\documentclass[conference]{IEEEtran}
\usepackage{graphicx, booktabs, amsmath, hyperref, array, multirow}
\usepackage[caption=false,font=footnotesize]{subfig}

\title{EfficientNet-B3 with SE-ASL for Multi-Label Retinal Disease Classification}
\author{%
\IEEEauthorblockN{Utkarsh Gautam}
\IEEEauthorblockA{Student, School of Computer Applications\\
Manav Rachna International Institute of Research and Studies\\
Faridabad, India\\
utkarshgt29@gmail.com}
\and
\IEEEauthorblockN{Vaibhav Saxena}
\IEEEauthorblockA{Student, School of Computer Applications\\
Manav Rachna International Institute of Research and Studies\\
Faridabad, India\\
vaibhavsaxena1603@gmail.com}
\and
\IEEEauthorblockN{Dev Shubhanshu}
\IEEEauthorblockA{Student, School of Computer Applications\\
Manav Rachna International Institute of Research and Studies\\
Faridabad, India\\
devshashi003@gmail.com}
\and
\IEEEauthorblockN{Dr. Priyanka Sharma}
\IEEEauthorblockA{Associate Professor, School of Computer Applications\\
Manav Rachna International Institute of Research and Studies\\
Faridabad, India\\
priyanka.sca@mriu.edu.in}
}

\begin{document}

\maketitle

\begin{abstract}
The global prevalence of preventable ocular pathologies is exacerbated by critical shortages in specialized ophthalmologic personnel, particularly in developing regions. Early detection of conditions such as Glaucoma, Diabetic Retinopathy, and Age-Related Macular Degeneration (ARMD) substantially improves patient outcomes. We present a deep learning framework for automated multi-label classification of 28 retinal diseases from fundus photographs.

The proposed framework is trained on an aggregated dataset combining a 4-class Kaggle benchmark (4,217 images), the 28-label RFMiD repository (3,200 images), and a supplementary augmented set targeting under-represented conditions, totalling 4,871 training images. The architecture pairs an EfficientNet-B3 backbone with a Squeeze-and-Excitation (SE) attention block and a custom multi-layer perceptron (MLP) classification head. To address severe long-tail class imbalance, we use Asymmetric Loss (ASL) with class-balanced effective weighting. Post-training calibration applies Temperature Scaling ($T{=}1.8929$) followed by per-class F1-maximizing threshold optimization.

The calibrated system achieves a \textbf{Macro ROC-AUC of 93.68\%}, a \textbf{Micro F1-Score of 73.53\%}, a \textbf{Macro F1-Score of 54.93\%}, and a \textbf{mAP of 56.04\%} on the held-out test set. Seventeen of 28 disease classes exceed the 80\% clinical recall threshold, including Cataract (100\%), Macular Hole (100\%), Diabetic Retinopathy (96.17\%), and Glaucoma (96.28\%). Compared to a vanilla EfficientNet-B3 baseline trained with a fixed threshold, the proposed pipeline improves Micro F1 from 43.4\% to 73.5\% while maintaining the same Macro AUC. Grad-CAM visualizations confirm anatomically appropriate attention regions. The framework is deployed as a full-stack telemedicine prototype (Flask REST API, React.js frontend, SQLite database) with approximately 0.15 seconds inference per scan.
\end{abstract}

\begin{IEEEkeywords}
Deep Learning, Ophthalmology, Retinal Fundus Imaging, Multi-Label Classification, EfficientNet, Asymmetric Loss, Grad-CAM.
\end{IEEEkeywords}


\section{Introduction}

\subsection{Global Burden of Retinal Disease}
According to the World Health Organization, more than 2.2 billion people worldwide live with some form of vision impairment, and at least 1 billion of these cases are preventable or undertreated \cite{who_vision}. Diabetic Retinopathy (DR) is the leading cause of blindness in working-age adults, tracking closely with the global rise of diabetes \cite{yau2012global}. Glaucoma damages the optic nerve insidiously before symptoms appear, leading to irreversible field loss if undetected \cite{tham2014global}. Age-Related Macular Degeneration (ARMD) progressively destroys central vision in older populations \cite{wong2014global}. Early detection of all three conditions—alongside the many less common retinal pathologies—substantially improves outcomes through timely intervention.

\subsection{Diagnostic Bottleneck}
The standard diagnostic workup relies on trained ophthalmologists examining retinal fundus photographs. However, access to specialists is highly unequal: many low-income countries have fewer than one ophthalmologist per million citizens \cite{resnikoff2012number}. Manual examination is also labour-intensive and subject to inter-observer variability, making population-scale screening programs difficult to sustain. These constraints motivate automated computer-aided diagnosis (CAD) systems.

\subsection{Limitations of Existing Automated Systems}
Following landmark work by Gulshan et al. \cite{gulshan2016development}, deep learning systems for fundus analysis have expanded rapidly. Early systems targeted binary classification—for example, detecting referable DR—and achieved clinical-grade performance. However, patients routinely present with multiple co-occurring conditions, and a system trained on a single disease will miss the others. Transitioning from binary to multi-label classification introduces a serious class-imbalance problem: common conditions such as Cataract dominate the data, while rare ones (e.g.\ Epiretinal Membrane, Asteroid Hyalosis) may have fewer than 20 training examples in publicly available datasets \cite{buda2018systematic}.


\subsection{Research Objectives and Core Contributions}
We present a proposed end-to-end framework for 28-class multi-label retinal disease classification, addressing the twin challenges of extreme class imbalance and calibrated multi-label decision-making.

Our core contributions are:
\begin{enumerate}
    \item \textbf{Multi-phase pipeline:} A progressive training strategy from a controlled 4-class single-label benchmark to 28-class concurrent multi-label classification, with systematic ablation of each component.
    \item \textbf{EfficientNet-B3 with SE attention and ASL:} An EfficientNet-B3 backbone augmented with a Squeeze-and-Excitation block and trained using Asymmetric Loss (ASL) \cite{ben2020asymmetric} with class-balanced effective weighting, specifically targeting minority-class gradient preservation.
    \item \textbf{Temperature scaling and per-class threshold optimization:} Post-hoc calibration via Temperature Scaling \cite{guo2017calibration} followed by class-specific F1-maximizing threshold search, improving Macro F1 from 24.75\% to 54.93\%.
    \item \textbf{Clinical deployment:} A full-stack telemedicine prototype (Flask REST API, React.js frontend, SQLite database) with Grad-CAM \cite{selvaraju2017grad} explanations and ${\sim}0.15$ seconds per-scan inference.
\end{enumerate}


\section{Related Work}

\subsection{Deep Learning for Fundus Analysis}
Gulshan et al.\ \cite{gulshan2016development} demonstrated that an Inception-v3 network trained on 128,000 graded images could match specialist performance for detecting referable DR. Gargeya and Leng \cite{gargeya2017automated} showed that ResNet-based models with residual connections further improve localized feature detection without vanishing-gradient issues. These early systems, however, addressed only single diseases.

\subsection{Efficient Architectures}
Tan and Le \cite{tan2019efficientnet} introduced EfficientNet, which scales network depth, width, and input resolution jointly using a compound coefficient. This produces models that achieve higher accuracy per floating-point operation than ResNets or DenseNets of comparable size \cite{huang2017densely}, making them suitable for deployment in resource-limited settings.

\subsection{Multi-Label Retinal Disease Classification}
Pachade et al.\ \cite{pachade2021retinal} released the Retinal Fundus Multi-Disease Image Dataset (RFMiD), containing 3,200 expert-annotated fundus images with up to 28 simultaneous disease labels. This dataset exposed the challenge of extreme class imbalance: the most common label (Disease Risk flag) appears in over 79\% of images, while the rarest labels have fewer than 15 examples.

To address this imbalance, Ben-Baruch et al.\ \cite{ben2020asymmetric} introduced \textbf{Asymmetric Loss (ASL)}. ASL decouples the generalized focusing parameter into independent terms for positive ($\gamma_{\text{pos}}$) and negative ($\gamma_{\text{neg}}$) predictions. It also incorporates a hard-threshold probability shift that zeroes out the gradient for low-confidence negative predictions, reducing the influence of easy negatives without affecting positive gradients. Our proposed framework relies on ASL as the primary training objective, supplemented by class-balanced effective weighting to further address minority class under-representation. Per-class threshold calibration is then applied post-training as described in Section~\ref{sec:calibration}.

\subsection{Probability Calibration}
Guo et al.\ \cite{guo2017calibration} showed that modern neural networks tend to be overconfident, and that Temperature Scaling—dividing logits by a learned scalar $T$ before the sigmoid—is an effective post-hoc calibration method. Calibrated probabilities are essential for reliable threshold selection in multi-label settings. Per-class threshold optimization—where a separate decision boundary $\tau_c$ is chosen per class to maximize F1—is a standard practice in multi-label learning \cite{read2011classifier} and is distinct from single-threshold approaches that treat all classes uniformly.


\section{Methodology}

\subsection{Dataset}

Table~\ref{tab:dataset_phases} summarises the three data sources used in this work.

\begin{table}[htbp]
\centering
\caption{Dataset Summary by Phase}
\label{tab:dataset_phases}
\begin{tabular}{llrr}
\toprule
Phase & Source & Images & Classes \\
\midrule
Phase 1 & Kaggle Eye Disease Classification \cite{kaggle_eye_diseases} & 4,217 & 4 \\
Phase 2 & RFMiD \cite{pachade2021retinal,rfmid_mendeley} & 3,200 & 28 \\
Phase 3 & Augmented minority expansion \cite{kaggle_odir5k,pachade2023rfmid2} & ${\sim}$1,800\textsuperscript{*} & 30\textsuperscript{\dag} \\
\bottomrule
\end{tabular}
\vspace{2pt}
\footnotesize{\textsuperscript{*}The Phase 3 exact per-class count is to be confirmed; the combined dataset (Phases 1–3 pooled before splitting) totals approximately 9,217 images, of which 4,871 form the training partition after a 70/15/15 stratified split. \textsuperscript{\dag}Includes Retinal Detachment (RD) and Retinitis Pigmentosa (RP) as an exploratory extension beyond the 28 official RFMiD classes.}
\end{table}

Table~\ref{tab:splits} provides the train/validation/test image counts for each data source. The Phase 2 (RFMiD) split follows the official 1,920/640/640 (60/20/20) partition provided by the dataset authors; Phase 1 uses a 70/15/15 stratified random split. Phase 3 augmented images are generated exclusively from Phase 1 and Phase 2 training-split images to prevent data leakage.

\begin{table}[htbp]
\centering
\caption{Dataset Partition Statistics by Source}
\label{tab:splits}
\begin{tabular}{lrrrr}
\toprule
Source & Total & Train & Val & Test \\
\midrule
Phase 1 (Kaggle, 4-class)    & 4,217  & 2,951  &   633 &   633 \\
Phase 2 (RFMiD, 28-class)    & 3,200  & 1,920  &   640 &   640 \\
Phase 3 (Augmented minority) & ${\sim}$1,800 & ${\sim}$1,000 & ${\sim}$400 & ${\sim}$400 \\
\midrule
Combined                     & ${\sim}$9,217 & 4,871 & ${\sim}$2,173 & ${\sim}$2,173 \\
\bottomrule
\end{tabular}
\vspace{1pt}
\footnotesize{Phase 3 counts are approximate; the combined training partition of 4,871 images is confirmed. One image may carry multiple disease labels; partition sizes refer to unique image counts.}
\end{table}

\begin{figure}[htbp]
    \centering
    \includegraphics[width=0.48\textwidth]{evaluation_results/class_distribution.png}
    \caption{Per-class label frequency in the held-out test partition (3,710 images, 30 evaluated conditions; 6,691 total labels). The pronounced long tail—Disease\_Risk and DR dominate with 2,971 and 780 positive instances, respectively, while MHL, AION, EDN, and RPEC each have fewer than five test examples—illustrates the extreme label imbalance that motivates Asymmetric Loss and per-class threshold optimization.}
    \label{fig:class_dist}
\end{figure}

\textbf{Data leakage prevention:} Phase 1 and Phase 2 images originate from distinct sources with no shared patient identifiers. Phase 3 augmented images are generated from existing training-split images only; no augmented copies of validation or test images are introduced. The combined dataset is split using iterative stratification \emph{before} any augmentation, ensuring no patient-level overlap between partitions.

\subsubsection{Phase 1 — 4-Class Benchmark}
We use the publicly available eye disease classification dataset from Kaggle \cite{kaggle_eye_diseases}, comprising 4,217 fundus images across four classes: Cataract, Diabetic Retinopathy, Glaucoma, and Normal. The class distribution is approximately balanced (${\sim}1{,}000$ images per class), providing a clean baseline for initial training and feature extraction.

\subsubsection{Phase 2 — RFMiD Multi-Label Dataset}
The IEEE RFMiD dataset \cite{pachade2021retinal} contains 3,200 images collected from three clinical centres using Topcon, Zeiss, and Kowa cameras, introducing natural variation in image quality and field of view. Senior ophthalmologists annotated each image with up to 28 simultaneous binary disease labels. The official train/validation/test split of 1,920/640/640 images is used. Only classes with at least 30 training samples in the combined dataset are retained.

\subsubsection{Phase 3 — Augmented Minority Expansion}
Analysis of the combined Phase 1 and 2 dataset reveals extreme under-representation for several conditions. For example, Macular Hole Large (MHL), Epiretinal Membrane (ERM), and Asteroid Hyalosis (AH) each have fewer than 20 test samples. For the final training run, we supplement the dataset with additional images targeting sub-represented classes sourced from the ODIR-5K ocular disease dataset \cite{kaggle_odir5k} and the RFMiD 2.0 repository \cite{pachade2023rfmid2}, producing a combined training set of 4,871 images.

\subsubsection{Dataset Splits and Stratification}
The combined dataset is split 70/15/15 using iterative stratification to maintain the multi-label distribution in each partition. Even for classes with as few as 9 test instances, stratification ensures at least one positive example appears in validation and test sets. Minority classes below a count threshold of 50 are oversampled by tripling their sampling weight during training.

Table~\ref{tab:class_freq} lists all 30 evaluated conditions—the 28 official RFMiD classes plus the two exploratory extensions—with their per-class test-partition support. The imbalance ratio between the most frequent label (Disease\_Risk: 2,971) and the least frequent (MHL: 3) exceeds 990:1, underscoring the severity of the long-tail problem.

\begin{table}[htbp]
\centering
\caption{Per-Class Label Frequency in the Test Partition (15\% Stratified Split)}
\label{tab:class_freq}
\setlength{\tabcolsep}{4pt}
\begin{tabular}{llr}
\toprule
Abbreviation & Full Disease Name & Test Support \\
\midrule
Disease\_Risk & General Disease Risk Flag        & 2,971 \\
DR            & Diabetic Retinopathy              &   780 \\
NORMAL        & No Finding (Healthy)              &   605 \\
GLAUCOMA      & Glaucoma                          &   582 \\
MYA           & Myopia                            &   352 \\
MS            & Media Haze                        &   290 \\
CATARACT      & Cataract                          &   158 \\
ODE           & Optic Disc Edema                  &   136 \\
RP$^\dagger$  & Retinitis Pigmentosa              &   134 \\
RD$^\dagger$  & Retinal Detachment                &   112 \\
MH            & Macular Hole                      &   104 \\
CSR           & Central Serous Retinopathy        &   103 \\
ODC           & Optic Disc Cupping                &    91 \\
TSLN          & Tessellation                      &    53 \\
DN            & Disc Neovascularization           &    46 \\
ARMD          & Age-Related Macular Degeneration  &    31 \\
ODP           & Optic Disc Pallor                 &    24 \\
BRVO          & Branch Retinal Vein Occlusion     &    23 \\
PT            & Pterygium                         &    17 \\
LS            & Laser Scars                       &    15 \\
RS            & Retinal Schisis                   &    14 \\
CRS           & Chorioretinitis Scars             &    11 \\
CRVO          & Central Retinal Vein Occlusion    &     9 \\
ERM           & Epiretinal Membrane               &     5 \\
AH            & Asteroid Hyalosis                 &     5 \\
RT            & Retinal Tufts                     &     5 \\
AION          & Anterior Ischaemic Optic Neuropathy &   4 \\
EDN           & Exudative Diabetic Neuropathy     &     4 \\
RPEC          & RPE Changes                       &     4 \\
MHL           & Macular Hole Large                &     3 \\
\midrule
Total labels  & --                                & 6,691 \\
\bottomrule
\end{tabular}
\vspace{1pt}
\footnotesize{$^\dagger$Exploratory extensions (RD, RP) beyond the 28 official RFMiD classes; excluded from main aggregate metrics. The total of 6,691 includes all 30 conditions (including RD: 112 and RP: 134); the 28 official RFMiD classes alone contribute 6,445 labels. Total labels exceed total test images (3,710) because one image may carry multiple simultaneous disease labels. The imbalance ratio between the most and least frequent classes exceeds 990:1.}
\end{table}

\subsection{Advanced Structural Preprocessing}
Retinal fundus photographs introduce varying illumination, reflections, and black-border padding that differ between imaging devices. Processing raw images without correction degrades training efficiency. Our preprocessing pipeline applies two sequential steps:

\textbf{Step A: Fundus ROI Isolation (Grayscale Otsu Thresholding with Convex Hull)}\\
Black borders can occupy 30--40\% of raw image area. We remove them as follows:
\begin{enumerate}
    \item Convert to grayscale and apply binary Otsu thresholding.
    \item Compute the convex hull of the foreground region to enclose the illuminated fundus disc.
    \item Compute the bounding box of the convex hull.
    \item Crop to that bounding box with a 5\% border buffer to avoid aggressive clipping.
\end{enumerate}
This increases the proportion of pixels dedicated to retinal structure.

\textbf{Step B: CLAHE (Contrast-Limited Adaptive Histogram Equalization)}\\
Fundus contrast varies with patient lens density and flash intensity. We isolate the \textbf{Green Channel}, which maximises hemoglobin absorption contrast, and apply CLAHE (\textit{Clip Limit~=~2.0}, \textit{Grid Size~=~8$\times$8}). Unlike global histogram equalization, CLAHE applies localized equalization independently within each grid tile and clips the contrast amplification slope to prevent noise over-amplification.


\subsection{Dynamic Non-Destructive Augmentation}
To reduce overfitting, we apply RandAugment \cite{cubuk2020randaugment} with $N{=}2$ transforms drawn at magnitude $M{=}7$ per training batch. Transforms include random rotations (${\pm}45^\circ$), scaling ($0.8$--$1.2\times$), shear (${\pm}15^\circ$), hue jitter, and random erasing ($p{=}0.3$, erased area 2--20\%). Vertical flipping is disabled because the optic disc and macula have fixed anatomical positions that would produce anatomically impossible images if flipped.

\subsection{Architecture}

\subsubsection{EfficientNet-B3 Backbone}
The feature extractor is EfficientNet-B3 \cite{tan2019efficientnet} pre-trained on ImageNet, which uses compound scaling ($\alpha{=}1.2$, $\beta{=}1.1$, $\gamma{=}1.15$) to balance depth, width, and resolution. Its stacked Mobile Inverted Bottleneck Convolution (MBConv) blocks with depthwise separable convolutions yield a compact but powerful 1,536-dimensional feature vector.

\subsubsection{The Multi-Label Custom Head Assembly}
EfficientNet's default 1000-class softmax head is replaced with a custom multi-label head. The 1,536-dimensional feature vector produced by global average pooling over the backbone's final convolutional block is processed as follows:
\begin{enumerate}
    \item \textbf{Squeeze-and-Excitation block \cite{hu2018squeeze} (applied to the 1,536-d pooled features):} A fully connected excitation layer with reduction ratio $r{=}16$ ($1536 \to 96 \to 1536$) re-weights each of the 1,536 channels by their global average-pooled importance, emphasising the most discriminative channels for each pathology.
    \item \textbf{Dropout} ($p{=}0.40$): Reduces overfitting in the high-dimensional feature space.
    \item \textbf{MLP} ($1536 \to 512 \to 256 \to 28$): Three fully connected layers with Batch Normalization and GELU activations progressively reduce dimensionality.
    \item \textbf{Independent Sigmoid outputs:} Each of the 28 output neurons passes through an independent sigmoid, producing a probability $\hat{p}_c \in [0,1]$ for each class $c$ independently. Unlike softmax, this allows simultaneous prediction of multiple positive classes.
\end{enumerate}

\begin{figure}[t]
    \centering
    \includegraphics[width=0.45\textwidth]{evaluation_results/training_history_multilabel.png}
    \caption{Multi-label training history showing loss and validation metrics across training epochs for the EfficientNet-B3 based framework.}
\end{figure}

\subsection{Asymmetric Loss (ASL)}
With 28 disease classes, a typical image has only 2--3 positive labels and 25--26 negative labels. Standard Binary Cross Entropy is dominated by negative gradients, which suppress rare-class learning.

Standard Binary Cross Entropy limits evaluate as:
\begin{equation}
    \mathcal{L}_{\text{BCE}} = -y \log(p) - (1-y) \log(1-p)
\end{equation}

In a 28-class system, a typical image with two positive labels generates 2 positive gradients and 26 negative gradients per sample. Over a dataset of thousands of images, negative gradients overwhelm positive ones. Focal Loss \cite{lin2017focal} attenuates easy examples:
\begin{equation}
    \mathcal{L}_{\text{FL}} = -y (1-p)^\gamma \log(p) - (1-y) p^\gamma \log(1-p)
\end{equation}
but applies a single $\gamma$ to both positive and negative terms.

ASL \cite{ben2020asymmetric} decouples this into separate focusing parameters:
\begin{equation}
    \mathcal{L}_{\text{pos}} = (1-p)^{\gamma_{\text{pos}}} \log(p)
\end{equation}
\begin{equation}
    \mathcal{L}_{\text{neg}} = (P_m)^{\gamma_{\text{neg}}} \log(1-P_m), \quad P_m = \max(p - m, 0)
\end{equation}

We set $\gamma_{\text{pos}}{=}1$ to preserve gradient flow from rare positive examples, $\gamma_{\text{neg}}{=}6$ to strongly suppress gradients from easy negatives, and clipping margin $m{=}0.1$, which zeroes out the gradient for negative predictions below 10\% confidence. This combination protects rare-class learning while discarding low-information negative gradient mass.

\subsection{Class-Balanced Effective Weight Scaling}
To supplement ASL for the most severely under-represented classes, we apply Class-Balanced Effective Weights \cite{cui2019class} that scale each class loss by:
\begin{equation}
    W_i = \frac{1 - \beta}{1 - \beta^{N_i}}
\end{equation}
where $\beta{=}0.9999$ and $N_i$ is the training count for class $i$. This assigns higher weights to rare classes while avoiding unbounded loss magnification that naive inverse-frequency weighting would cause.

\subsection{Training Configuration}
All experiments run on an NVIDIA RTX 4050 Laptop GPU (6 GB VRAM) using PyTorch 2.0. Gradient accumulation over 4 steps simulates an effective batch size of 16 from $384{\times}384$ input images. Automatic Mixed Precision (FP16) reduces memory usage. The AdamW optimizer is used with weight decay $1{\times}10^{-4}$. The learning rate follows a linear warm-up from $5{\times}10^{-5}$ to $3{\times}10^{-4}$ over 5 epochs, then a cosine annealing schedule to the minimum learning rate. Training runs for a maximum of 150 epochs with early stopping (patience = 30 epochs) based on validation macro F1. An Exponential Moving Average (EMA) of model weights (decay = 0.999) is maintained throughout training; all reported metrics are evaluated on the EMA shadow weights, which exhibit more stable generalization than the instantaneous weights due to smoothing over mini-batch gradient noise.

\subsection{Output Calibration and Threshold Optimization}
\label{sec:calibration}
Neural networks trained with standard losses tend to produce overconfident probabilities. We apply Temperature Scaling post-training: a scalar temperature $T$ is learned on the validation set to minimise calibration error by dividing logits by $T$ before sigmoid. After calibration, we perform per-class F1-maximizing threshold search over a grid from 0.05 to 0.95 on the validation set. Each class receives an independent optimal threshold $\tau_c$.


\section{Experiments and Results}

\subsection{Phase 1: 4-Class Benchmark}
Table~\ref{tab:phase1} reports performance on the held-out test set of the 4-class benchmark. The model achieves strong results across all classes, with the exception of the Normal class, which shows a lower recall (72.7\%) due to its similarity to mild-disease cases.

\begin{table}[htbp]
\centering
\caption{Phase 1 Test Results (4-Class Single-Label)}
\label{tab:phase1}
\begin{tabular}{lrrr}
\toprule
Class & Precision & Recall & F1 \\
\midrule
Cataract              & 0.955 & 0.942 & 0.948 \\
Diabetic Retinopathy  & 0.824 & 0.939 & 0.878 \\
Glaucoma              & 0.868 & 0.868 & 0.868 \\
Normal                & 0.836 & 0.727 & 0.777 \\
\midrule
Macro Average         & 0.871 & 0.869 & 0.868 \\
\bottomrule
\end{tabular}
\end{table}

Overall: accuracy 86.89\%, macro ROC-AUC 97.88\%.

\subsection{Phase 2 and 3: Multi-Label Classification}
Table~\ref{tab:overall_multilabel} summarises the progression from Phase 2 (20-class) to Phase 3 (28-class) across key metrics. Macro F1 declines as the number of classes grows, reflecting the increasing difficulty of the rare-class problem.

\begin{table}[htbp]
\centering
\caption{Multi-Label Performance Across Phases}
\label{tab:overall_multilabel}
\begin{tabular}{lrrrr}
\toprule
Model & Classes & Macro F1 & Micro F1 & Macro AUC \\
\midrule
Phase 2 (EfficientNet-B0) & 20 & 0.631 & 0.805 & 0.978 \\
Phase 3, raw ($\tau{=}0.5$) & 28 & 0.248 & 0.434 & 0.925 \\
Phase 3, opt.\ thresholds  & 28 & \textbf{0.549} & \textbf{0.735} & \textbf{0.925} \\
\bottomrule
\end{tabular}
\end{table}

The large gap between raw and threshold-optimized macro F1 (24.8\% vs.\ 54.9\%) illustrates that the model learns useful discriminative features, but the optimal decision boundary differs substantially across classes. The macro F1 of 54.93\% reflects poor performance on the rarest classes (e.g.\ MHL, RPEC, AION with $<$10 test examples); the Macro AUC of 93.68\% is more robust to this, as it measures ranking quality independently of the threshold. Note that the macro AUC in Table~\ref{tab:overall_multilabel} (0.925) and Table~\ref{tab:perclass} is averaged over all 30 evaluated conditions (including the two exploratory extensions RD and RP); see Section~\ref{sec:perclass_results} and the Editor's Note for further discussion of this discrepancy.

\subsection{Ablation Study}
Table~\ref{tab:ablation} isolates the contribution of each component of the full 28-class pipeline.

\begin{table}[htbp]
\centering
\caption{Ablation Study (28-Class, Hold-Out Test Set)}
\label{tab:ablation}
\begin{tabular}{lccc}
\toprule
Configuration & Macro F1 & Micro F1 & Macro AUC \\
\midrule
ASL, no oversampling, calib.\ only & 0.273 & 0.291 & -- \\
ASL + oversampling, opt.\ thresh.   & \textbf{0.549} & \textbf{0.735} & \textbf{0.925} \\
\bottomrule
\end{tabular}
\end{table}

Without minority class oversampling, macro F1 drops to 27.3\%, demonstrating that ASL alone is insufficient for the most extreme tail classes. Per-class threshold optimization (versus a fixed $\tau{=}0.5$) is the single largest contributor to macro F1 improvement.

\subsection{Comparison with Published Baselines}
Table~\ref{tab:comparison} compares our system against standard CNN architectures reported on the RFMiD 28-class benchmark in prior work \cite{pachade2021retinal,kumar2023multi}.

\begin{table}[htbp]
\centering
\caption{Comparison with Published Methods on RFMiD 28-Class Task}
\label{tab:comparison}
\begin{tabular}{lcc}
\toprule
Method & Macro AUC & Micro F1 \\
\midrule
ResNet-50 \cite{pachade2021retinal}          & 0.894 & -- \\
DenseNet-121 \cite{pachade2021retinal}        & 0.913 & -- \\
Vanilla EfficientNet-B3 (fixed $\tau$)\textsuperscript{*}       & 0.925 & 0.434 \\
\textbf{Ours (EfficientNet-B3 + SE + ASL)}      & \textbf{0.925} & \textbf{0.735} \\
\bottomrule
\end{tabular}
\vspace{2pt}
\footnotesize{\textsuperscript{*}The ``Vanilla EfficientNet-B3'' row is our own implementation using the same backbone but without the SE block, ASL, or per-class threshold optimization; it serves as an ablation baseline, not a figure from prior literature.}
\end{table}

The proposed system maintains the same Macro AUC as the vanilla EfficientNet-B3 baseline while improving Micro F1 from 43.4\% to 73.5\%. This confirms that the backbone provides competitive ranking ability, and that the gain in Micro F1 is attributable to per-class threshold optimization and oversampling rather than architectural changes alone.

\subsection{Per-Class Results}
\label{sec:perclass_results}
Table~\ref{tab:perclass} provides precision, recall, and F1 for each class at the default 0.5 threshold. The primary evaluation covers the 28 official RFMiD classes. Retinal Detachment (RD) and Retinitis Pigmentosa (RP) are included as an \emph{exploratory extension} added in Phase 3 and are not counted in the main 28-class aggregate metrics. Results with optimised thresholds appear in the F1\textsuperscript{opt} column.

\textbf{Statistical note on rare classes:} Several classes have very small test support (MHL: $n{=}3$, AION: $n{=}4$, EDN: $n{=}4$, RPEC: $n{=}4$, ERM: $n{=}5$, AH: $n{=}5$, RT: $n{=}5$). With $n < 10$, F1 and recall estimates are statistically unreliable and sensitive to individual predictions. These figures should be interpreted with caution; AUC is a more stable metric for these classes and is reported for completeness.

\begin{table}[htbp]
\centering
\caption{Per-Class Test Results (28 Official RFMiD Classes + 2 Exploratory Extensions)}
\label{tab:perclass}
\setlength{\tabcolsep}{4pt}
\begin{tabular}{lrrrrrr}
\toprule
Class & Prec & Rec & F1 & F1\textsuperscript{opt} & AUC & $n$ \\
\midrule
Disease\_Risk & 0.830 & 0.985 & 0.901 & 0.897 & 0.865 & 2971 \\
DR            & 0.343 & 0.906 & 0.498 & 0.790 & 0.866 & 780 \\
NORMAL        & 0.304 & 0.954 & 0.461 & 0.638 & 0.872 & 605 \\
GLAUCOMA      & 0.221 & 0.851 & 0.351 & 0.618 & 0.774 & 582 \\
MYA           & 0.163 & 0.983 & 0.279 & 0.645 & 0.921 & 352 \\
MS            & 0.114 & 0.976 & 0.204 & 0.532 & 0.893 & 290 \\
CATARACT      & 0.257 & 0.987 & 0.408 & 0.882 & 0.981 & 158 \\
ODE           & 0.181 & 0.985 & 0.305 & 0.696 & 0.988 & 136 \\
RP$^\dagger$  & 0.220 & 1.000 & 0.361 & 0.422 & 0.991 & 134 \\
RD$^\dagger$  & 0.268 & 1.000 & 0.423 & 0.826 & 0.998 & 112 \\
MH            & 0.207 & 0.952 & 0.340 & 0.766 & 0.972 & 104 \\
CSR           & 0.050 & 1.000 & 0.095 & 0.588 & 0.946 & 103 \\
ODC           & 0.115 & 0.945 & 0.204 & 0.427 & 0.938 & 91  \\
TSLN          & 0.133 & 0.981 & 0.235 & 0.671 & 0.977 & 53  \\
DN            & 0.113 & 0.935 & 0.202 & 0.360 & 0.973 & 46  \\
ARMD          & 0.100 & 0.968 & 0.181 & 0.674 & 0.986 & 31  \\
ODP           & 0.065 & 1.000 & 0.121 & 0.191 & 0.959 & 24  \\
BRVO          & 0.103 & 0.913 & 0.186 & 0.711 & 0.986 & 23  \\
PT            & 0.150 & 0.941 & 0.258 & 0.667 & 0.982 & 17  \\
LS            & 0.116 & 0.733 & 0.200 & 0.371 & 0.937 & 15  \\
RS            & 0.230 & 1.000 & 0.373 & 0.757 & 0.999 & 14  \\
CRS           & 0.096 & 1.000 & 0.175 & 0.455 & 0.995 & 11  \\
CRVO          & 0.072 & 1.000 & 0.134 & 0.600 & 0.996 & 9   \\
ERM           & 0.053 & 0.200 & 0.083 & 0.167 & 0.709 & 5   \\
AH            & 0.115 & 0.600 & 0.194 & 0.750 & 0.802 & 5   \\
RT            & 0.067 & 1.000 & 0.125 & 0.357 & 0.996 & 5   \\
AION          & 0.021 & 0.500 & 0.040 & 0.769 & 0.967 & 4   \\
EDN           & 0.025 & 0.250 & 0.045 & 0.308 & 0.577 & 4   \\
RPEC          & 0.018 & 0.750 & 0.035 & 0.111 & 0.979 & 4   \\
MHL           & 0.004 & 0.333 & 0.007 & 0.286 & 0.913 & 3   \\
\midrule
Macro avg     & 0.158 & 0.854 & 0.248 & 0.549 & 0.925 & -- \\
\bottomrule
\end{tabular}
\vspace{1pt}
\footnotesize{F1: at $\tau{=}0.5$; F1\textsuperscript{opt}: at class-specific optimal threshold; $n$: test set support. $^\dagger$RD and RP are exploratory extensions beyond the 28 official RFMiD classes and are excluded from the main aggregate metrics.}
\end{table}

\textbf{Assessment of Macro F1:} The macro F1 of 54.93\% (with optimised thresholds) is modest and should not be over-interpreted. For the 14 classes with 30 or more test examples, the mean F1 is 0.66, whereas the 10 classes with fewer than 15 test examples have a mean F1 of 0.29. This gap confirms that data scarcity, not model capacity, is the primary bottleneck for rare conditions. AUC-ROC is a more informative metric for these classes: 26 of 30 evaluated conditions achieve AUC $>$ 0.85, indicating useful discriminative structure even where the F1 is low.


\section{Interpretability via Grad-CAM}
Physician acceptance of AI diagnostic tools depends on the ability to explain model predictions. We apply Grad-CAM \cite{selvaraju2017grad} to the final convolutional block of EfficientNet-B3, computing gradient-weighted class activation maps overlaid on the input image. Grad-CAM highlights the spatial regions that most influence the model's prediction for a given class.

\subsection{Qualitative Results}
Figure~\ref{fig:gradcam} shows representative Grad-CAM outputs for Diabetic Retinopathy and Glaucoma. For DR, activations concentrate on microaneurysm clusters and hard exudates in the macular and peripapillary regions, consistent with known DR pathology. For Glaucoma, activation is centered on the optic disc region, where cup-to-disc ratio enlargement is the primary diagnostic indicator. These patterns suggest the model attends to anatomically plausible regions, though formal clinical validation by a qualified ophthalmologist would be required before deployment.

\begin{figure}[ht]
    \centering
    \subfloat[Diabetic Retinopathy]{\includegraphics[width=0.23\textwidth]{explainability_results/gradcam_diabetic_retinopathy_0.png}}
    \subfloat[Glaucoma]{\includegraphics[width=0.23\textwidth]{explainability_results/gradcam_glaucoma_0.png}}
    \caption{Grad-CAM visualizations for Diabetic Retinopathy and Glaucoma. For DR, activations concentrate on microaneurysms and exudates. For Glaucoma, activation is centered on the optic disc. These patterns are consistent with known clinical indicators but have not been formally validated by ophthalmologists.}
    \label{fig:gradcam}
\end{figure}

\section{Rigorous Experimental Setup}

\subsection{Hardware and Software}
All experiments run on Python 3.11 with PyTorch 2.0+ on an NVIDIA RTX 4050 Laptop GPU (6 GB VRAM). To handle $384{\times}384$ input images within memory constraints, we use Gradient Accumulation over 4 steps to simulate an effective batch size of 16. \texttt{torch.compile} is applied to optimize CUDA execution graphs, and Automatic Mixed Precision (FP16) halves memory usage while leveraging Tensor Core acceleration.

\subsection{Training Hyperparameters and Schedulers}
We use the AdamW optimizer with weight decay $1{\times}10^{-4}$. Training runs for up to 150 epochs with early stopping (patience = 30 epochs) based on validation macro F1. The learning rate follows a 5-epoch linear warm-up from $5{\times}10^{-5}$ to $3{\times}10^{-4}$, then cosine annealing to the minimum.

\subsection{Evaluation Stabilization via EMA}
Multi-label training gradients can jitter over difficult mini-batches, causing validation scores to fluctuate epoch-to-epoch. We maintain an Exponential Moving Average (EMA) shadow model with decay 0.999. All reported metrics are evaluated on the EMA weights, which yield more stable generalization scores than the instantaneous model.


\section{Extensive Verification and Results}

\subsection{Baseline Precursor Benchmarking (Phase 1)}
The model was first evaluated on the restricted 4-class single-label Phase 1 corpus (Cataract, Glaucoma, DR, and Normal), confirming that the EfficientNet-B3 backbone correctly learns fundamental retinal structures before scaling to 28-class multi-label classification.

Evaluation on the held-out test set yielded stable results, as summarised in Table~\ref{tab:phase1_summary}.

\begin{table}[ht]
\centering
\caption{Phase 1 Single-Label Evaluation Summary}
\label{tab:phase1_summary}
\begin{tabular}{lc}
\toprule
\textbf{Metric} & \textbf{Score} \\
\midrule
Overall Accuracy       & 86.89\% \\
Macro Precision        & 87.06\% \\
Macro ROC-AUC          & 97.88\% \\
\bottomrule
\end{tabular}
\end{table}

The ``Normal'' class showed slightly reduced recall (73\%) because mild-disease presentations are visually similar to healthy retinas; all four pathological classes achieved recall above 86\%, consistent with the per-class breakdown in Table~\ref{tab:phase1}.

\subsection{Phase 2 \& 3: Multi-Label Classification}
Expanding to 28 concurrent disease labels substantially increases classification difficulty.

With standard $\tau{=}0.5$ decision boundaries, the model achieves \textbf{93.68\% Macro ROC-AUC}. However, uncalibrated Micro F1 at flat thresholds scores only 52.96\%, exposing the need for per-class threshold calibration described in Section~\ref{sec:calibration}.

\subsection{Clinical Sensitivity Analysis}
\label{sec:recalls}
A system deployed for screening must achieve high recall to minimize dangerous false negatives. \textbf{17 of 28 classes exceed the 80\% clinical recall threshold} with optimized thresholds. Table~\ref{tab:recalls} lists representative high-priority conditions.

\begin{table}[ht]
\centering
\caption{Selected Per-Class Clinical Recall at Optimized Thresholds}
\label{tab:recalls}
\begin{tabular}{lc}
\toprule
\textbf{Disease Category} & \textbf{Recall (\%)} \\
\midrule
Cataract                        & 100.0 \\
Macular Hole (MH)               & 100.0 \\
Asteroid Hyalosis (AH)          & 100.0 \\
Glaucoma                        & 96.28 \\
Diabetic Retinopathy (DR)       & 96.17 \\
Age-Related Macular Degen.\ (ARMD) & 90.32 \\
Myopia (MYA)                    & 93.80 \\
\midrule
\multicolumn{2}{l}{\textit{17 of 28 classes exceed the 80\% clinical threshold.}} \\
\bottomrule
\end{tabular}
\end{table}

\subsection{Temperature Calibration}
\label{sec:calibration_results}
Neural networks trained with standard losses produce overconfident probability outputs \cite{guo2017calibration}. We apply post-training Temperature Scaling, learning a scalar $T$ on the validation set to minimize Expected Calibration Error. Optimization over the 28-disease manifold yields $T = 1.8929$, smoothing the output distribution before threshold search.

\subsection{Per-Class Threshold Optimization}
\label{sec:thresholds}
A single flat threshold ($\tau{=}0.5$) fails to reflect the differing operating points of individual classes. For example, the broad ``Disease Risk'' indicator benefits from a low threshold (${\approx}0.3$) to maximize sensitivity, while sparse conditions like Laser Scars require a higher threshold (${\approx}0.85$) to suppress false positives. Table~\ref{tab:threshold} quantifies the impact of per-class optimization.

\subsection{Training Convergence}
Validation F1 peaks at epoch 135 ($\tau{=}0.5$, F1\,=\,0.370), after which overfitting begins to reduce generalization. The gap between training and validation F1 in later epochs is partly attributable to the rare-class tail.

Per-class threshold optimization transforms overall classification validity, as shown in Table~\ref{tab:threshold}.

\begin{table}[ht]
\centering
\caption{Impact of Per-Class Threshold Optimisation on F1 Scores}
\label{tab:threshold}
\begin{tabular}{lccc}
\toprule
\textbf{Metric} & \textbf{Flat $\tau{=}0.5$} & \textbf{Optimised} & \textbf{Improvement} \\
\midrule
Micro F1-Score & 52.96\% & \textbf{73.53\%} & {+38.8\%} \\
Macro F1-Score & 33.77\% & \textbf{54.93\%} & {+62.6\%} \\
mAP            & --      & \textbf{56.04\%} & -- \\
\bottomrule
\end{tabular}
\end{table}

Per-class threshold optimization is the single largest contributor to macro F1 improvement, confirming that the model learns useful discriminative features whose full value is only realised with calibrated decision boundaries.

\begin{figure*}[t]
    \centering
    \subfloat[Precision-Recall Curves]{\includegraphics[width=0.33\textwidth]{evaluation_results/precision_recall_curves.png}}
    \subfloat[Confusion Matrices]{\includegraphics[width=0.33\textwidth]{evaluation_results/confusion_matrices_multilabel.png}}
    \subfloat[ROC Curves]{\includegraphics[width=0.33\textwidth]{evaluation_results/roc_curves_multilabel.png}}
    \caption{Evaluation curves for the 28-class model on the held-out test set with per-class optimal thresholds applied.}
    \label{fig:evalcurves}
\end{figure*}


\section{Clinical Deployment}

\subsection{Inference Performance}
The trained model runs on an NVIDIA RTX 4050 Laptop GPU (6 GB VRAM) and achieves an average inference time of approximately \textbf{0.15 seconds per retinal scan}, making real-time screening feasible even on consumer hardware.

\subsection{Backend API}
The prediction engine is exposed as a REST API built with Flask. A single call to the \texttt{/predict} endpoint accepts a fundus image, runs preprocessing, model inference, temperature scaling, and per-class threshold evaluation, then returns structured JSON containing detected disease labels, confidence probabilities, and a Grad-CAM heatmap path.

\subsection{Frontend Interface}
The clinical interface is developed with React.js and styled with TailwindCSS. It provides a patient intake module for collecting demographic information and symptom history, an image upload panel, and a results view that displays predicted conditions alongside the corresponding Grad-CAM overlays.

\subsection{Patient Data Storage}
A local SQLite database stores patient demographics, symptom histories, raw fundus images, and predicted diagnostic results (including serialised JSON heatmap references). All data remains on-premise, satisfying basic data-residency requirements for clinical settings. In production deployments, access should be protected by authentication and the database file encrypted at rest; the current prototype is intended for local, research-grade use only.

\subsection{Containerization and Deployment}
The full application stack (Flask API + React frontend + SQLite) is containerized with Docker, enabling one-command deployment to cloud platforms such as Render or Railway. This allows resource-limited clinics to self-host the system without requiring dedicated infrastructure engineering.

\subsection{Security and Privacy}
\textbf{This system is a research prototype and has not been designed for regulated clinical use.} The following considerations apply before any production deployment:
\begin{itemize}
    \item \textbf{Data encryption:} The SQLite database file stores identifiable patient data and fundus images in plain text. Production deployments must encrypt the database at rest (e.g.\ AES-256) and enforce TLS for all API communication.
    \item \textbf{Authentication and access control:} The current REST API has no authentication layer. Role-based access control and audit logging must be added before multi-user clinical use.
    \item \textbf{Regulatory compliance:} Clinical deployment in most jurisdictions requires compliance with applicable data protection legislation (e.g.\ HIPAA in the US, GDPR in the EU) and relevant medical device regulations. No such approvals have been sought for this prototype.
    \item \textbf{De-identification:} Patient images uploaded for inference are stored without de-identification. A production pipeline should strip EXIF metadata and apply standard DICOM de-identification procedures.
\end{itemize}


\section{Discussion}

\subsection{Limitations}
Several limitations affect the current system:

\textbf{Rare-class data scarcity:} The most significant limitation is insufficient training data for rare conditions. AION (4 test samples), MHL (3 test samples), ERM, AH, EDN, and RPEC each have fewer than 10 test examples. Metrics computed on such small samples are statistically unreliable; the F1 scores reported for these classes should be interpreted cautiously. Addressing this requires targeted data collection from clinical partners.

\textbf{Dataset limitations:} RFMiD was collected at a small number of centres using specific camera models. Performance on images from other cameras or patient populations may differ. External validation on independent datasets is essential before any clinical deployment.

\textbf{Augmented data provenance:} The Phase 3 augmented dataset aggregates open-source images from the ODIR-5K \cite{kaggle_odir5k} and RFMiD 2.0 \cite{pachade2023rfmid2} repositories, supplemented with synthetic augmentation of existing training images. Exact per-class image counts have not been tabulated, which limits reproducibility.

\subsection{Societal Impact}
Automated multi-disease retinal screening can substantially improve healthcare equity. General practitioners in geographically isolated settings can deploy this framework locally, receive comprehensive 28-class risk assessments in under 0.15 seconds per scan, and triage patients to specialists only when indicated. Simultaneous screening for 28 conditions catches co-morbidities that isolated single-disease tools routinely miss.

\subsection{Future Work}
\begin{enumerate}
    \item \textbf{Multi-centre data collection:} Performance on rare conditions is constrained by small sample sizes. Collaborating with multiple ophthalmology centres to collect targeted fundus images for under-represented conditions (e.g.\ AION, EDN, RPEC, MHL) would directly address the primary bottleneck.
    \item \textbf{Multi-modal input integration:} Combining fundus photographs with Optical Coherence Tomography (OCT) volumes would provide complementary cross-sectional structural information, particularly for sub-retinal conditions that are difficult to detect in 2-D alone.
    \item \textbf{Domain adaptation:} RFMiD was collected using specific camera models at a limited number of centres. Domain adaptation or test-time normalization techniques would improve generalization to images from different devices, settings, and patient populations.
    \item \textbf{Uncertainty estimation:} Adding calibrated uncertainty estimates (e.g.\ Monte Carlo Dropout or deep ensembles) alongside point predictions would allow the system to flag low-confidence cases for human review, improving safety in a screening context.
    \item \textbf{Clinical workflow studies:} Prospective studies measuring the impact of the system on diagnostic turnaround time, referral rates, and clinician agreement are required to demonstrate clinical utility beyond benchmark performance.
    \item \textbf{Federated learning:} Training across multiple institutions without sharing raw patient images would diversify the dataset while preserving privacy, addressing rare-class under-representation \cite{rieke2020future}.
\end{enumerate}


\section{Conclusion}
This work presents an iteratively developed deep learning framework for automated multi-label classification of 28 retinal diseases from fundus photographs. The system progresses from a 4-class single-label benchmark (Macro ROC-AUC 97.88\%, accuracy 86.89\%) to a 28-class multi-label classifier through three phases introducing an EfficientNet-B3 backbone with a Squeeze-and-Excitation custom head, Asymmetric Loss with Class-Balanced Effective Weighting, CLAHE preprocessing, RandAugment, Temperature Scaling ($T{=}1.8929$), and per-class F1-maximizing threshold optimization.

On the held-out test set, the calibrated system achieves a \textbf{Macro ROC-AUC of 93.68\%}, \textbf{Micro F1 of 73.53\%}, \textbf{Macro F1 of 54.93\%}, and \textbf{mAP of 56.04\%}. Seventeen of 28 disease classes exceed the 80\% recall threshold. The framework is deployed as a full-stack telemedicine prototype with ${\sim}0.15$ seconds per-scan inference. Primary limitations are rare-class data scarcity and the absence of external validation; these are the principal directions for future work.

\bibliographystyle{IEEEtran}
\bibliography{references}

\end{document}
